\section{Einleitung}

Quaternionwertige neuronale Netzwerke (QVNNs) verarbeiten Eingaben, Gewichte und Aktivierungen nicht als einzelne reelle Zahlen, sondern als Quaternionen. Ein Quaternion
\[
q = q_0 + q_1\mathbf{i} + q_2\mathbf{j} + q_3\mathbf{k}
\]
besteht aus vier gekoppelten reellen Komponenten. Das ist fuer Daten interessant, bei denen mehrere Komponenten zusammengehoeren, etwa bei Rotationen, Farbinformationen oder mehrdimensionalen Signalen. In solchen Faellen koennen QVNNs kompaktere Modelle und teilweise bessere Generalisierung liefern als vergleichbare reellwertige Netzwerke \cite{bill2023comparison,article}. Gleichzeitig ist die Optimierung schwieriger, weil das Hamilton-Produkt nicht kommutativ ist. Die Reihenfolge von Faktoren ist also relevant; Standard-Backpropagation laesst sich nicht ohne Anpassung uebertragen.

Aus Sicht des Moduls OML ist genau diese Optimierungsfrage zentral. Bill, Cox und Champagne vergleichen mehrere Quaternion-Backpropagation-Verfahren und zeigen, dass bereits kleine Unterschiede in Lernregel, Produktreihenfolge und Hyperparametern deutliche Auswirkungen auf Genauigkeit und Laufzeit haben koennen \cite{bill2023comparison}. In ihrer Studie erreicht der Clifford Multilayer Perceptron (CMLP) auf mehreren Regressionsbenchmarks die staerkste Gesamtleistung. Dieses Projekt nimmt CMLP deshalb als Optimierungsbaseline und fragt, ob Predictive Coding eine brauchbare Alternative oder Erweiterung fuer QVNNs sein kann.

Predictive Coding Networks (PCNs) verwenden keine reine Feedforward-Aktivierungskette. Stattdessen versucht jede Schicht, die Aktivitaet der darunterliegenden Schicht vorherzusagen. Bei festen Gewichten werden zuerst latente Zustaende iterativ inferiert; erst danach werden die Gewichte anhand lokaler Vorhersagefehler aktualisiert \cite{stenlund2025pcn}. Ein PCN ist deshalb nicht nur ein anderer Optimierer fuer ein MLP, sondern eine eigene Lernarchitektur mit zwei Zeitskalen: schnelle Inferenz und langsameres Gewichtslernen. Dadurch entstehen zusaetzliche Optimierungsentscheidungen, insbesondere Inferenzschrittweite, Anzahl Inferenzschritte, latente Initialisierung, Lernrate und Batchgroesse.

Die Fragestellung des Projekts lautet daher:

\begin{quote}
Laesst sich ein quaternionwertiges neuronales Netzwerk mit Predictive Coding auf Standardaufgaben zur Approximation nichtlinearer Funktionen stabil optimieren, und wie verhaelt sich diese Optimierung im Vergleich zu starken Quaternion-Backpropagation-Baselines?
\end{quote}

Der Bericht konzentriert sich auf diese Optimierungsfrage. Die theoretischen Grundlagen werden nur so weit eingefuehrt, wie sie fuer Methodik und Interpretation notwendig sind. Der Schwerpunkt liegt auf den implementierten Modellen, den Hyperparametern, den Resultaten und den beobachteten Trainingsdynamiken.
